% ============================================================
%  بخش سه: مقاله‌ی تحلیلی–تطبیقی (متن آکادمیک اصلی)
% ============================================================
\section{تحلیل تطبیقی مفهوم آزادی در فلسفه‌ی سیاسی}
\label{sec:main}

% ─────────────────────────────────────────────────────────────
\subsection{مقدمه و روش‌شناسی}
\label{subsec:intro}
% ─────────────────────────────────────────────────────────────

مفهوم آزادی یکی از پُرمناقشه‌ترین واژگان در فلسفه‌ی سیاسی است. به گفته‌ی \thinker{آیزایا برلین}، آزادی آنقدر معانی گوناگون دارد که «تاریخ این کلمه، آینه‌ی تاریخ بشریت است.»\footnote{\lr{Berlin, I. (1969). \textit{Four Essays on Liberty}. Oxford University Press, p. 121.}} این مقاله تلاش می‌کند نه آزادی را تعریف کند، بلکه \textbf{جغرافیای تعریف‌ها} را ترسیم نماید: این تعریف‌ها در پاسخ به چه پرسش‌هایی ظهور کردند، با چه رقیبانی درگیر شدند، و چه پیامدهایی داشتند. % FIXED: \lr in footnote

روش‌شناسی این مقاله ترکیبی از سه رویکرد است:

\begin{enumerate}
  \item \textbf{تحلیل تاریخی-مفهومی} (\lr{Cambridge School — Skinner, Pocock}): هر مفهوم در بستر زبانی و تاریخی خود فهمیده می‌شود؛ % FIXED: already has \lr
  \item \textbf{تحلیل تطبیقی ساختاری}: مقایسه‌ی نظام‌مند تعابیر مختلف بر اساس پنج معیار: دشمن مفهومی، انسان‌شناسی، رابطه‌ی فرد/جامعه، موضع درباره‌ی دولت، و قلمرو کاربرد؛
  \item \textbf{تحلیل تنش‌های درونی}: شناسایی تناقضات و تنش‌های ناگزیر در هر نظریه.
\end{enumerate}

% ─────────────────────────────────────────────────────────────
\subsection{آزادی منفی: از هابز تا برلین}
\label{subsec:negative}
% ─────────────────────────────────────────────────────────────

\subsubsection{صورت‌بندی اولیه: هابز}

\thinker{توماس هابز} در \textit{لویاتان} (۱۶۵۱) آزادی را به‌صورتی کاملاً مکانیکی تعریف کرد: «یک انسان آزاد کسی است که در انجام آنچه اراده، قدرت، و فهم دارد، هیچ مانعی نداشته باشد.»\footnote{\lr{Hobbes, T. (1651/1994). \textit{Leviathan}. Indianapolis: Hackett, ch. 21.}} این تعریف عمداً فرد را \textit{اتم فیزیکی} فرض می‌کند که در فضای اجتماعی حرکت می‌کند: آزادی=غیبت مانع خارجی. پیامد مهم این است که قانون—هر قانونی—به خودی خود محدودیت آزادی است، حتی اگر توسط خود فرد پذیرفته شده باشد. % FIXED: \lr in footnote

این صورت‌بندی بنیادین ضعفی دارد: اگر کسی در قفسِ طلایی باشد که هیچ‌گاه در آن توقفی نمی‌کند (چون هرگز نمی‌خواهد از قفس خارج شود)، آیا آزاد است؟ برلین این مشکل را بعدها با مثال «بیماری» توضیح داد.

\subsubsection{کمال‌یابی با آیزایا برلین}

مقاله‌ی کلاسیک \thinker{آیزایا برلین}، «دو مفهوم از آزادی» (۱۹۵۸)، تمایز تأثیرگذاری میان \concept{آزادی منفی} و \concept{آزادی مثبت} کشید:

\begin{tcolorbox}[defbox, title={تمایز برلین}]
\begin{multicols}{2}
\textbf{\color{PrimaryDeep}آزادی منفی (\lr{Negative}):}\\ % FIXED: \lr
«آزادی \textit{از}»\\
پرسش: آیا کسی مانع من است؟\\
دشمن: اجبار و مداخله\\
مقیاس: حوزه‌ی بدون مانع\\

\columnbreak

\textbf{\color{EraContemp}آزادی مثبت (\lr{Positive}):}\\ % FIXED: \lr
«آزادی \textit{برای}»\\
پرسش: چه کسی ارباب من است؟\\
دشمن: وابستگی، عقلانیت ناقص\\
مقیاس: کنترل بر خود و زندگی\\
\end{multicols}
\end{tcolorbox}

برلین نگران آن بود که آزادی مثبت می‌تواند به \enemy{توتالیتاریسم} منجر شود: اگر دولت ادعا کند که «خودِ واقعی» شما را—که ظاهراً می‌خواهد در سوسیالیسم زندگی کند—از «خود تجربی» شما که می‌خواهد سیگار بکشد، جدا کند، دولت می‌تواند اجبار را در نقاب آزادی پنهان کند.

\subsubsection{نقدهای وارد بر برلین}

\thinker{چارلز تیلور} استدلال کرد که تمایز برلین ساده‌انگارانه است: آزادی منفی هم همیشه به ارزیابی‌ای درباره‌ی اهمیت اهداف نیازمند است. ما همیشه بین موانع مهم و کم‌اهمیت تمایز می‌گذاریم—یعنی ناخودآگاه به آزادی مثبت تکیه داریم.\footnote{\lr{Taylor, C. (1979). What's wrong with negative liberty? In A. Ryan (Ed.), \textit{The Idea of Freedom}. Oxford.}} % FIXED: \lr in footnote

\thinker{جرالد مکالوم} پیشنهاد کرد که آزادی همیشه یک رابطه‌ی سه‌تایی است: \textit{X از Y آزاد است برای انجام Z}. هر اظهاری درباره‌ی آزادی ضمناً درباره‌ی هر سه پارامتر حرف می‌زند.\footnote{\lr{MacCallum, G. (1967). Negative and positive freedom. \textit{Philosophical Review}, 76(3), 312--334.}} % FIXED: \lr in footnote

% ─────────────────────────────────────────────────────────────
\subsection{آزادی مثبت: از روسو تا تیلور}
\label{subsec:positive}
% ─────────────────────────────────────────────────────────────

سنت آزادی مثبت ریشه در ایده‌ی خودفرمانی دارد. \thinker{ژان ژاک روسو} معتقد بود آزادی واقعی نه اطاعت از قانون دیگری، بلکه اطاعت از قانونی است که خود انسان وضع کرده: «اطاعت از قانونی که خود برای خود نوشته‌ایم آزادی است.»\footnote{\lr{Rousseau, J.-J. (1762/1968). \textit{The Social Contract}. Harmondsworth: Penguin, Book I, ch. 8.}} % FIXED: \lr in footnote

این سنت از طریق \thinker{کانت}—که خودقانون‌گذاری عقل عملی را کانون اخلاق و آزادی دانست—به ایدئالیسم آلمانی و سپس به لیبرالیسم تجدیدنظرطلب انگلیسی (\thinker{تی.اچ. گرین، بوزانکه}) رسید.

\thinker{تی.اچ. گرین} مهم‌ترین صورت‌بندی را ارائه داد: آزادی واقعی نه «رهایی از همه‌ی محدودیت‌ها»، بلکه «قدرت مثبت برای انجام یا لذت‌بردن از چیزی ارزشمند» است.\footnote{\lr{Green, T.H. (1881). Liberal legislation and freedom of contract. In R.L. Nettleship (Ed.), \textit{Works of T.H. Green}. London.}} این صورت‌بندی آزادی را با \textit{توسعه‌ی توانایی‌های انسانی} پیوند زد و پایه‌ی دولت رفاه را ریخت. % FIXED: \lr in footnote

% ─────────────────────────────────────────────────────────────
\subsection{آزادی جمهوری‌خواهانه: عدم سلطه}
\label{subsec:republican}
% ─────────────────────────────────────────────────────────────

\thinker{فیلیپ پتیت} در \textit{جمهوری‌خواهی} (۱۹۹۷) رویکرد سومی معرفی کرد که ادعا می‌کند از هر دو سنت قبلی رادیکال‌تر است:

\begin{tcolorbox}[defbox, title={آزادی به مثابه عدم سلطه}]
آزادی منفی برلین می‌گوید: آزادم مادامی که کسی مداخله نکند. اما این کافی نیست. یک برده در دست ارباب مهربانی که اتفاقاً هرگز مداخله نمی‌کند همچنان \textbf{برده} است—چون ارباب \textit{می‌تواند} مداخله کند. \textbf{عدم سلطه} یعنی دیگران حتی \textit{قدرت} مداخله‌ی دلبخواهانه را نداشته باشند.\footnote{\lr{Pettit, P. (1997). \textit{Republicanism: A Theory of Freedom and Government}. Oxford: Clarendon Press, p. 52.}} % FIXED: \lr in footnote
\end{tcolorbox}

پتیت برای تأمین عدم سلطه، به نهادهای جمهوری‌خواه نیاز دارد: حاکمیت قانون، تفکیک قوا، مشارکت مدنی، و مقاومت برابر تجمع قدرت. این دیدگاه در دو دهه‌ی اخیر در نظریه‌ی دموکراسی مشارکتی و نقد نئولیبرالیسم بسیار تأثیرگذار بوده است.

\thinker{کوئنتین اسکینر} از منظر تاریخ‌نگاری کمبریج استدلال کرد که این مفهوم اصالتاً ریشه‌ی جمهوری‌خواه رومی دارد—در برابر سنت هابزی که آن را حذف کرد.\footnote{\lr{Skinner, Q. (1998). \textit{Liberty before Liberalism}. Cambridge University Press.}} % FIXED: \lr in footnote

% ─────────────────────────────────────────────────────────────
\subsection{رویکرد قابلیت: سن و نوسبام}
\label{subsec:capability}
% ─────────────────────────────────────────────────────────────

\thinker{آمارتیا سن} رویکرد قابلیت (\lr{Capability Approach}) را در اواخر دهه‌ی ۱۹۸۰ معرفی کرد که بعداً \thinker{مارتا نوسبام} آن را به شکل جامع‌تری توسعه داد. پرسش محوری این رویکرد: آزادی فرمال کافی نیست؛ مهم است که آیا مردم واقعاً \textit{قادر} به استفاده از آزادی‌های خود هستند؟

\begin{tcolorbox}[defbox, title={قابلیت‌های مرکزی نوسبام}]
\small
نوسبام فهرستی از ده قابلیت اساسی ارائه می‌دهد که هر جامعه‌ای باید آنها را تضمین کند، از جمله: زندگی، سلامت جسمی، تمامیت جسمانی، حواس/تخیل/اندیشه، احساسات، عقل عملی، ارتباط اجتماعی، ارتباط با دنیای طبیعی، بازی، و کنترل محیط سیاسی.\footnote{\lr{Nussbaum, M. (2011). \textit{Creating Capabilities: The Human Development Approach}. Harvard University Press, pp. 33--34.}} % FIXED: \lr in footnote
\end{tcolorbox}

\concept{آزادی در رویکرد قابلیت} دوگانه است: \textit{آزادی فرصت} (داشتن گزینه‌ها) و \textit{آزادی فرآیند} (تعیین سرنوشت). هر دو ضروری هستند. این رویکرد نزدیک‌ترین نظریه به سنت آزادی مثبت است، اما با پایه‌ی تجربی‌تر و حساسیت بیشتر به تفاوت‌های فردی و بستری.

% ─────────────────────────────────────────────────────────────
\subsection{آزادی در مکاتب راست، چپ، و میانه}
\label{subsec:leftright}
% ─────────────────────────────────────────────────────────────

یکی از جالب‌ترین ابعاد مفهوم آزادی این است که هر طیف سیاسی ادعا می‌کند حامی اصیل آن است. جدول زیر این مواضع را به‌صورت تطبیقی نشان می‌دهد:

\begin{table}[h!]
\centering
\caption{آزادی در مکاتب سیاسی: مقایسه‌ی تطبیقی}
\label{tab:leftright}
\renewcommand{\arraystretch}{1.6}
\setlength{\tabcolsep}{5pt}

\begin{tabularx}{\textwidth}{|>{\bfseries}p{2.4cm}
  |>{\raggedright}X|>{\raggedright}X|>{\raggedright}X
  |>{\raggedright\arraybackslash}X|}
\hline
\rowcolor{PrimaryDeep}
\textcolor{white}{\textbf{معیار}} &
\textcolor{white}{\textbf{راست کلاسیک\\(لیبرتاریانیسم)}} &
\textcolor{white}{\textbf{لیبرالیسم میانه}} &
\textcolor{white}{\textbf{سوسیال‌دموکراسی}} &
\textcolor{white}{\textbf{مارکسیسم}} \\
\hline

\rowcolor{BgCool}
تعریف آزادی &
عدم مداخله‌ی دولت &
عدم مداخله + امکانات &
قابلیت‌های واقعی &
رهایی از کار بیگانه \\
\hline

\rowcolor{white}
دشمن اصلی &
دولت و مالیات &
استبداد + فقر &
نابرابری ساختاری &
سرمایه‌داری \\
\hline

\rowcolor{BgCool}
رابطه با دولت &
دولت حداقلی &
دولت میانجی &
دولت رفاه &
دولت انتقالی \\
\hline

\rowcolor{white}
انسان‌شناسی &
فرد اتمیستیک &
فرد اجتماعی &
شهروند اجتماعی &
موجود طبقاتی \\
\hline

\rowcolor{BgCool}
مشروعیت نابرابری &
کامل (بازار) &
مشروط (رالز) &
محدود &
نفی کامل \\
\hline

\rowcolor{white}
آزادی اقتصادی &
اولویت اول &
با بازتوزیع &
تنظیم‌شده &
جمعی‌شده \\
\hline

\rowcolor{BgCool}
متفکران شاخص &
هایک، فریدمن، نوزیک &
رالز، دورکین، میل &
برلین (اصلاح‌شده)، سن &
مارکس، لوکاچ، گرامشی \\
\hline

\end{tabularx}
\end{table}

\subsubsection{تنش‌های ساختاری: سه بحران اصلی}

\paragraph{۱. تنش فرد/جامعه:}
همه‌ی نظریه‌های آزادی با این چالش دست‌وپنجه نرم می‌کنند که آزادی فرد تا کجاست و آزادی جمع از کجا آغاز می‌شود. لیبرالیسم کلاسیک به فرد تقدم می‌دهد؛ جمهوری‌خواهی به فضای عمومی؛ مارکسیسم به طبقه.

\paragraph{۲. تنش صوری/واقعی:}
آزادی فرمال (حق رأی برابر) و آزادی واقعی (قدرت استفاده از آن حق) همیشه یکی نیستند. رویکرد قابلیت و سوسیال‌دموکراسی این شکاف را جدی‌ترین معضل می‌دانند.

\paragraph{۳. تنش آزادی/امنیت:}
هر گسترش آزادی‌های خاصی، گاهی امنیت یا برابری را به خطر می‌اندازد. این تنش در فلسفه‌ی سیاسی همیشه «حل‌نشده» باقی می‌ماند.

% ─────────────────────────────────────────────────────────────
\subsection{آزادی اگزیستانسیالیستی: انتخاب، مسئولیت، هستی}
\label{subsec:exist}
% ─────────────────────────────────────────────────────────────

سنت اگزیستانسیالیستی آزادی را نه مسئله‌ی حقوق یا نهادها، بلکه \textit{ساختار هستی انسان} می‌داند. \thinker{ژان‌پل سارتر} معروف‌ترین گزاره را بیان کرد: «وجود بر ماهیت تقدم دارد» — انسان چیزی نیست مگر آنچه خود می‌سازد؛ از این‌رو «محکوم به آزادی» است.\footnote{\lr{Sartre, J.-P. (1946/2007). \textit{Existentialism is a Humanism}. New Haven: Yale University Press, p. 29.}} % FIXED: \lr in footnote

\thinker{هانا آرنت} از زاویه‌ای سیاسی‌تر، آزادی را با \concept{کنش} (\lr{action}) پیوند زد: آزادی تنها در فضای عمومی و در لحظه‌ی \textit{شروع امر نو} تحقق می‌یابد. این نگاه عمیقاً ضد روان‌شناسانه است—آزادی در درون ذهن نیست، بلکه در «بین ما» (\lr{inter homines esse}) است.\footnote{\lr{Arendt, H. (1961). \textit{Between Past and Future}. New York: Viking, pp. 143--171.}} % FIXED: \lr in footnote and text

% ─────────────────────────────────────────────────────────────
\subsection{آزادی فمینیستی و پسااستعماری}
\label{subsec:fem}
% ─────────────────────────────────────────────────────────────

فمینیسم به‌خصوص از دهه‌ی ۱۹۷۰، نشان داد که مفاهیم لیبرالی آزادی بر انسانی \textit{نرینه، سفیدپوست، و طبقه‌ی متوسط} بنا شده‌اند. \thinker{کارول پیتمن} استدلال کرد که قرارداد اجتماعی اصلاً یک «قرارداد جنسی» است که زنان را در حوزه‌ی خصوصی محبوس می‌کند.\footnote{\lr{Pateman, C. (1988). \textit{The Sexual Contract}. Stanford University Press.}} % FIXED: \lr in footnote

سنت پسااستعماری (\thinker{فانون، اسپیواک، بابا}) نشان داد که مفهوم مدرن آزادی با استعمار هم‌زمان بوده: اروپا درباره‌ی آزادی انسان سخن می‌گفت، در حالی که بخش عظیمی از بشریت را استثمار می‌کرد. آزادی پسااستعماری باید این تاریخ تراژیک را در خود حمل کند.

\newpage
